\chapter{Conclusion}
\label{ch:conclusion}


\section{Summary of Findings}

This thesis asked how much of the reported performance of real-time sepsis
prediction on MIMIC-IV comes from the choice of model, and how much from the
researcher's choices about the label, the timing of prediction relative to
treatment, the data split and the evaluation metric. Across the six evaluated
model specifications, model choice affected AUROC more than the three Sepsis-3
operationalisations did. External transport has the second-largest point
estimate, but its interval is too wide to establish that rank, and the label
and the split moved AUROC least. The main effect of label choice was on cohort
membership and onset timing rather than on AUROC. The six pre-registered
hypotheses were resolved as follows.

\begin{enumerate}
    \item \textbf{H1 was not confirmed.} The label spread was
    \pcSpreadLabel{} AUROC against \pcSpreadModelClass{} across the six model
    specifications, and the interval on the difference,
    [\pcHOneCiLo, \pcHOneCiHi], excludes zero. Label choice nevertheless
    changed cohort membership substantially: the two narrowest readings agree
    at only $\kappa = \pcKappaVsBA$ despite near-identical prevalence.

    \item \textbf{H2 was met descriptively.} \pcHTwoEst{} of the
    \pcHTwoNTerms{} shared covariates met the pre-registered instability
    criterion, but only \pcHTwoNRobust{} of the changes, all in indicators of
    whether a test was ordered, exceeded their estimated noise.

    \item \textbf{H3 was not evaluable as pre-registered.} The pre-treatment
    event count was fixed at zero by the onset rule
    (Proposition~\ref{prop:h3}); standard re-timing (Variant F, post-hoc)
    recovered \pcPreTreatOnsetsVarF{} pre-treatment onsets on the same cohort.

    \item \textbf{H4 was not confirmed.} Random splitting did not produce
    higher AUROC: the random-minus-temporal difference was \pcHFourEst{}
    (95\% CI \pcHFourCiLo--\pcHFourCiHi), and every model fitted under both
    splits scored lower on the random split, although the two test sets also
    differ in composition.

    \item \textbf{H5 was not evaluable as pre-registered.} The internal Utility
    Score was too close to zero for a proportional decline to be defined, and
    the substituted external AUROC contrast, \pcHFiveEst{}
    (95\% CI \pcHFiveCiLo--\pcHFiveCiHi) on \pcExtNEventsB{} events, was
    imprecise. Every external point estimate was below its internal
    counterpart, and the post-hoc anchor-independent analysis showed a
    transport loss for the deterioration target, but the Sepsis-3 arm
    establishes neither the existence nor the magnitude of a degradation.

    \item \textbf{H6 supported one-sided non-inferiority, not equivalence.} GBT
    did not outperform the primary model by the 0.02 margin: the difference
    was \pcHSixEst{} [\pcHSixCiLo, \pcHSixCiHi], with the lower limit outside
    the margin. The results do not support preferring GBT over the primary
    model on predictive performance, and the small AUROC difference was not
    accompanied by a meaningful utility difference (\pcUtilityMaxGbtB{}
    against \pcUtilityMaxPrimaryB{} at each model's best threshold).
\end{enumerate}


\section{The Central Lesson}

The main consequence of label choice was not a large change in AUROC but a
change in who was labelled septic and when onset was assigned. Under the
pre-registered timing rule, pre-treatment prediction was structurally
impossible; standard re-timing recovered pre-treatment events on the same
cohort. Across all pre-registered labels, no evaluated operating point
approached the alert-burden benchmark: the best configuration under the
primary label reached a Utility Score of \pcUtilityMaxGbtB{} at
\pcUtilityOptBurdenGbtB{} false alerts per true alert, against a benchmark of
1.4, while AUROC remained in the high 0.7s. These findings indicate that
improving discrimination alone cannot resolve limitations introduced by outcome
membership, onset timing and alert burden. They do not isolate the label as the
sole cause of the gap between discrimination and usable alerting: the
estimand, the threshold, calibration, prevalence and the Utility Score's own
penalties all bear on it.

Splitting the Sepsis-3 conjunction (Section~\ref{sec:label_anchor_results})
both qualifies and sharpens this. Under every pre-registered variant the onset
time is fixed by the treatment anchor, so part of the variation between
variants follows from the definition. An anchor-independent deterioration
label agrees with Sepsis-3 at $\kappa = \pcKappaVsBD$, close to chance, so the
disagreement is not confined to readings of Sepsis-3. The treatment anchor
alone is about as predictable as the conjunction built on it
($\phi_{\text{anchor}} = \pcPhiAnchorPrimary$
[\pcPhiAnchorCiLoPrimary, \pcPhiAnchorCiHiPrimary] for the primary model, an
interval including 1, and $\pcPhiAnchorGbt$
[\pcPhiAnchorCiLoGbt, \pcPhiAnchorCiHiGbt] for GBT), whereas Sepsis-3 under
standard onset timing produced higher discrimination for the primary model
(\pcAltAurocFullPrimaryF{} against \pcAurocPrimaryB), on an in-horizon event
set slightly larger than Variant B's. Reduced discrimination is therefore
associated more with an onset placed at the clinician's action than with the
conjunction itself.


\section{Contributions Revisited}

\begin{enumerate}
    \item \textbf{C1 (Label variation, separated from discrimination):} Three
    Sepsis-3 operationalisations were applied to a discrete-time
    competing-risks model on \pcNStays{} MIMIC-IV stays, with the label and
    specification spreads estimated on the same test rows, each with an
    interval; to the best of our knowledge this combination is not previously
    reported. Label choice moved cohort membership substantially while moving
    AUROC by \pcSpreadLabel.

    \item \textbf{C2 (Treatment leakage: membership and onset timing):} The
    study proves that the conjunction-based onset rule leaves the
    pre-treatment window empty, and measures \pcPreTreatOnsetsVarF{}
    pre-treatment onsets on the same cohort under standard timing. Kamran et
    al.'s \cite{kamran2024evaluation} concern therefore holds for label
    membership, while for onset timing it depends on the onset rule.

    \item \textbf{C3 (Utility ceiling):} The Utility Score was swept across the
    full threshold grid for every model and label. Under the primary label its
    ceiling lies within a few per cent of the no-alert strategy and is reached
    only at \pcUtilityOptBurdenGbtB{} false alerts per true alert, extending
    Wang et al.'s \cite{wang2025srpm} meta-analytic finding to a controlled
    single-study comparison that does not depend on one operating point.

    \item \textbf{C4 (Coefficient instability):} \pcHTwoEst{} of the
    \pcHTwoNTerms{} shared covariates were unstable across labels under the
    pre-registered criterion; referred to their cluster-robust standard
    errors, \pcHTwoNRobust{} remain, all order indicators. This supports
    Lauritsen et al.'s \cite{lauritsen2021framing} framing-instability concern
    for specific effect estimates on MIMIC-IV.

    \item \textbf{C5 (Equity by label):} Subgroup discrimination and subgroup
    label sensitivity were measured on the same strata of one onset-prediction
    cohort under false-discovery control (Section~\ref{sec:search_strategy}).
    The label-sensitivity contrasts survive correction and the discrimination
    contrasts do not, and the surviving racial contrasts are missingness
    categories, so the finding concerns record completeness rather than race.

    \item \textbf{C6 (Reproducible pipeline):} A reproducible, seed-fixed
    pipeline built around a pre-registered protocol, with every amendment
    disclosed and reporting following TRIPOD+AI. Within the pinned environment
    it reproduces the reported estimates from the raw databases, and it writes
    the numbers in this document as generated constants
    (Section~\ref{sec:code_availability}, Appendix~\ref{app:code}).

    \item \textbf{C7 (Definitional vs.\ empirical label variation
    separated):} An anchor-independent deterioration label agrees with
    Sepsis-3 at $\kappa = \pcKappaVsBD$ and matches its onset hour in
    \pcPctExactMatchD\% of shared stays, and records \pcPreTreatOnsetsVarD{}
    onsets in the pre-treatment window that is empty of Sepsis-3 onsets under
    the pre-registered rule. These results show that label membership and
    onset timing shape the prediction task; they do not show that the
    physiology is easy to predict, since Variant D's discrimination level is
    set in part by its construction.
\end{enumerate}


\section{Future Work}

\begin{enumerate}
    \item \textbf{Better sepsis labels.} Future work should evaluate outcome
    definitions that reduce direct dependence on treatment timestamps and are
    adjudicated independently of the model inputs, for example definitions
    built on lactate trajectories and organ-function biomarkers. The
    deterioration label of Section~\ref{sec:label_anchor_results} is an
    operational construct, not a validated definition.

    \item \textbf{Lagged predictors.} A design that lags every predictor by one
    hour, or applies a sub-hour timestamp cut-off, would show whether the
    information associated with the recorded onset hour supports actionable
    advance warning (Limitation~\ref{lim:current_hour}).

    \item \textbf{Non-Sepsis-3 consensus definitions.} SEP-1 and the CDC Adult
    Sepsis Event were not implemented. Both are defined on antibiotic days, so
    neither tests treatment-independence, but each yields a different cohort,
    and Dutta et al.\ \cite{dutta2026performance} report that performance
    varies between them.

    \item \textbf{Prospective validation.} A new label should be tested where
    sepsis onset can be determined by a protocol-blinded assessor,
    independently of treatment decisions.

    \item \textbf{Utility-optimised training.} The models here are fitted to a
    likelihood or surrogate loss and only afterwards read at the
    utility-maximising threshold. Training against the Utility Score, through a
    differentiable surrogate or cost-sensitive reweighting, would change what
    the model learns. Such optimisation should use clinically justified costs
    and validation thresholds defined independently of the test set; otherwise
    it may improve the selected utility score without improving
    deployability.

    \item \textbf{Multi-hospital evaluation.} Whether any model can be
    deployed at an acceptable alert burden requires prospective multi-centre
    testing with clinical-staff feedback.
\end{enumerate}


\section{Closing Remark}

The central hypothesis was not supported: across the evaluated
specifications, model choice affected AUROC more than the three Sepsis-3
operationalisations did. The smaller AUROC effect should not, however, be
interpreted as evidence that label choice is unimportant. The variants selected
substantially different cohorts, and the onset rule determined whether a
pre-treatment prediction task could be defined at all.

No evaluated configuration approached the pre-specified alert-burden
benchmark. In addition, the reported discrimination concerns detection of the
recorded onset hour and does not establish advance warning. Future studies
should therefore report label agreement, onset timing, calibration, clinical
utility and alert burden alongside discrimination.

The principal contribution of this study is to show that outcome construction
changes the population and timing of the prediction task in ways that AUROC
alone does not reveal.
